\chapter{Conclusiones y Recomendaciones}
\section{Conclusiones}
Una vez planteado el problema general, definidos los objetivos e hipótesis
correspondientes, construido el sistema de inteligencia electoral Electoral-Peru-2026 y
aplicada la metodología CRISP-DM sobre un corpus multicanal del ecosistema
digital electoral peruano del año 2026, se exponen a continuación las
conclusiones que se desprenden de cada objetivo del estudio.
\subsection*{Respecto al Objetivo General}
El estudio cumplió el objetivo general al haber diseñado y puesto en marcha un
sistema de inteligencia electoral apoyado en Modelos de Lenguaje de Gran Escala
(LLMs), minería de sentimientos y detección dinámica de temas, aplicado a un
ecosistema digital multi-fuente del caso peruano. La plataforma Electoral-Peru-2026 articula
cuatro canales de información —portales de noticias, publicaciones en X/Twitter,
transcripciones de YouTube y publicaciones en TikTok— y procesa el corpus diario
íntegro (655 documentos válidos correspondientes al 22 de marzo de 2026) en
torno a 57 minutos de ejecución secuencial, con un costo aproximado de
US\$0{,}13 y un \textit{cache hit rate} del 85{,}5\,\% sobre el \textit{prompt}
de extracción. A partir de ese procesamiento, el sistema produce indicadores de
la percepción proyectada por los medios —sentimiento por candidato, temas
dominantes y distribución de cobertura por medio— accesibles en lenguaje
natural español mediante un servidor MCP que expone 14 herramientas
tipadas. La calidad de salida del sistema fue evaluada cuantitativamente
contra dos \textit{gold datasets} construidos para esta tesis: un dataset
de etiquetado con 100 documentos anotados por humano para los módulos de
sentimiento y temas, y un \textit{golden set} de 30 preguntas anotadas con
\textit{ground truth} de herramientas esperadas y respuesta de referencia
para la capa GraphRAG, comparada contra un RAG \textit{baseline} por
similitud semántica.
\subsection*{Respecto al Primer Objetivo Específico (OE1)}
Se diseñó e implementó la \textbf{arquitectura de ingesta y armonización
multi-fuente del sistema} sobre el patrón Medallion en tres capas de
procesamiento (Bronze, Silver y Gold), desplegado sobre una infraestructura
dual de bases de datos —PostgreSQL con \texttt{pgvector} para datos
relacionales y vectoriales, y Neo4j 5.20 Community Edition para el grafo de
conocimiento—. La arquitectura integró efectivamente las \textbf{cuatro
fuentes digitales} declaradas en el alcance (portales de noticias,
X/Twitter, YouTube y TikTok) bajo un esquema canónico común, preservando la
trazabilidad documento$\rightarrow$segmento$\rightarrow$análisis y la
calidad semántica del corpus. Como evidencia cuantitativa del cumplimiento,
sobre el corpus del 22 de marzo de 2026 la cascada de filtros pre-LLM
parametrizada desde catálogos YAML versionados
(\texttt{filter\_by\_account}, \texttt{filter\_by\_hashtag} y
\texttt{filter\_by\_title\_keywords}) redujo el volumen procesado de
\textbf{818 a 655 documentos válidos} ($\sim$20\,\% de descarte),
con un ahorro asociado en costo de \textit{tokens} de aproximadamente
\textbf{40\,\%}; la idempotencia del pipeline se verificó mediante tres
corridas consecutivas con conteos constantes en Postgres
(\texttt{documentos}, \texttt{segmentos}, \texttt{analisis\_llm}) y en Neo4j
(\texttt{Documento}, \texttt{Fragmento}, \texttt{Tema}). El OE1 se considera
\textbf{plenamente alcanzado}.

\subsection*{Respecto al Segundo Objetivo Específico (OE2)}
Se implementó el \textbf{modelo de análisis semántico contextual basado en
LLMs} para el etiquetado de sentimiento dirigido (ABSA) y la detección
dinámica de tópicos electorales emergentes, y se midió su desempeño tanto
operacionalmente como contra \textit{ground truth} humano. La configuración
final del módulo ABSA —\textbf{DeepSeek-Chat} en modalidad \textit{few-shot}
con 4 ejemplos calibrados para el español político peruano— operó con
\textbf{validez de schema del 100\,\%} (655/655 documentos parseables contra
el modelo Pydantic \texttt{AnalisisLLM} en el primer intento, sin reintentos),
\textbf{cobertura del 100\,\%} sobre el catálogo cerrado de 17 categorías
(sin alucinación de etiquetas) y un \textit{cache hit rate} medido de
\textbf{85,5\,\%} que reduce el costo operacional efectivo en
aproximadamente 75\,\%. Sobre el \textit{gold dataset} de 100 documentos,
el módulo obtuvo \textbf{F1 macro 0.856} en categoría temática (exactitud
0.830 sobre 17 clases), \textbf{F1 0.974} en relevancia geográfica al Perú,
\textbf{F1 macro 0.916} en identificación de candidato protagonista y
\textbf{MAE 0.107} con correlación Pearson de 0.892 en sentimiento continuo,
superando los umbrales declarados en la hipótesis específica HE2
(F1~$\geq$~0,80 y MAE~$\leq$~0,15). En la tarea de detección dinámica de
temas, el enfoque híbrido LLM~+~\textit{retrieval} multinivel detectó
\textbf{27 temas únicos} en el corpus del día, con \textbf{tasa de reuso}
del catálogo de 70,5\,\% en la corrida completa y \textbf{exactitud
semántica de 0.700} en la asignación documento$\rightarrow$tema sobre el
mismo \textit{gold dataset}. El único componente con desempeño limitado fue
la extracción multi-etiqueta de actores principales (F1 micro 0.295), cuya
naturaleza se interpreta como un sub-objetivo del \textit{prompt} y no como
un fallo del modelo. El OE2 se considera \textbf{alcanzado con métricas
superiores a los umbrales declarados}.

\subsection*{Respecto al Tercer Objetivo Específico (OE3)}
Se construyó el \textbf{grafo de conocimiento electoral} en Neo4j Community
Edition (versión 5.20), denominado \textbf{Electoral-Peru-2026}, mediante
extracción semántica del corpus etiquetado por el módulo LLM. El grafo
quedó conformado por \textbf{10 tipos de nodos} (\texttt{Documento},
\texttt{Fragmento}, \texttt{Candidato}, \texttt{Tema}, \texttt{Categoria},
\texttt{Medio}, \texttt{Actor}, \texttt{Region}, \texttt{Periodo} y
\texttt{Encuesta}) y \textbf{7 tipos de relaciones tipadas}
(\texttt{MENCIONA}, \texttt{SOBRE\_TEMA}, \texttt{PUBLICADO\_POR},
\texttt{PERTENECE\_A}, entre otras), articulando candidatos, temas, medios,
fragmentos y períodos. Sobre el subgrafo del día 22 de marzo de 2026 se
poblaron \textbf{723 nodos \texttt{Fragmento}}, \textbf{27 nodos
\texttt{Tema}} y \textbf{13 nodos \texttt{Categoria}}, con propiedades
\texttt{sentimiento\_doc} y \texttt{tema} materializadas sobre las aristas
\texttt{MENCIONA} y \texttt{SOBRE\_TEMA} respectivamente. La calidad
relacional del grafo se verificó mediante: (i) la correspondencia 1:1 entre
los conteos de \texttt{Fragmento}/\texttt{Documento} en Neo4j y los conteos
de \texttt{segmentos}/\texttt{documentos} en Postgres (idempotencia
cross-store); (ii) la ausencia de aristas huérfanas o nodos desconectados en
las consultas de auditoría; y (iii) la \textbf{ejecución correcta en una
única \textit{query} estructural Cypher} de las 14 herramientas del servidor
MCP que articulan las relaciones definidas, validando empíricamente la
hipótesis HE3. El OE3 se considera \textbf{alcanzado}.

\subsection*{Respecto al Cuarto Objetivo Específico (OE4)}
Se implementó el sistema \textbf{GraphRAG Electoral} sobre el grafo Neo4j
construido en el OE3, bajo una arquitectura basada en \textbf{herramientas
tipadas} (\textit{typed tools}) sobre el protocolo \textbf{MCP} (\textit{Model
Context Protocol}), en lugar del paradigma clásico de generación dinámica de
Cypher mediante \textit{few-shot prompting}. El servidor expone
\textbf{14 herramientas} agrupadas en seis dominios (candidatos, temas,
actores, evidencia, panorama, encuestas y macroeconómico), cada una con
esquemas Pydantic v2 de entrada y salida y consultas SQL/Cypher pre-cargadas
y parametrizadas. La validación operacional \textit{end-to-end} mostró que
las 14 herramientas retornaron \texttt{structuredContent} válido contra sus
esquemas Pydantic v2 en el \textbf{100\,\%} de las llamadas, con
\textbf{latencia mediana} entre 5 y 60\,ms para consultas no vectoriales y
de aproximadamente 1100\,ms para búsquedas semánticas, frente a los
1--2\,segundos típicos de un sistema GraphRAG con generación dinámica de
\textit{queries}. Sobre el \textit{golden set} de 30 preguntas anotadas, el
LLM cliente seleccionó el conjunto completo de herramientas esperadas en el
\textbf{83,3\,\%} de las preguntas y al menos una herramienta esperada en
el \textbf{100\,\%}, combinando recuperación semántica textual (búsquedas
sobre \texttt{pgvector}) con razonamiento estructural sobre el grafo
(\texttt{MATCH}/\texttt{MERGE} en Cypher) para producir respuestas
fundamentadas en evidencia textual del corpus original. El OE4 se considera
\textbf{alcanzado}.

\subsection*{Respecto al Quinto Objetivo Específico (OE5)}
Se evaluó el desempeño técnico y analítico del sistema integrado mediante
métricas semánticas y estructurales, contrastando las salidas del sistema
contra \textbf{dos \textit{gold datasets}} construidos específicamente para
esta tesis. (i) El \textit{gold dataset} de etiquetado (n=100 documentos,
validación asistida por dos evaluadores humanos) reportó \textbf{F1 macro
0.856} en categoría temática, \textbf{F1 0.974} en relevancia geográfica,
\textbf{F1 macro 0.916} en candidato protagonista, \textbf{MAE 0.107} en
sentimiento continuo y \textbf{exactitud semántica 0.700} en la asignación
documento$\rightarrow$tema. (ii) El \textit{golden set} de consulta (n=30
preguntas con \textit{ground truth} de herramientas esperadas y respuesta
de referencia) reportó \textbf{tool selection accuracy de 83,3\,\%},
\textbf{corrección de respuesta de 0.792} (frente a 0.196 del RAG
\textit{baseline} por similitud vectorial sobre el mismo corpus,
$\Delta=+0.596$), \textbf{relevancia de 0.842} (vs.\ 0.388,
$\Delta=+0.454$), \textbf{fidelidad de 0.854} y \textbf{similitud de
\textit{embeddings} de 0.721} (vs.\ 0.579, $\Delta=+0.142$), evaluadas por
un juez LLM independiente (DeepSeek-Chat, $T=0.2$). Todas las métricas
relevantes superaron los umbrales declarados en las hipótesis específicas
HE2 y HE4 y en la hipótesis general. El OE5 se considera \textbf{alcanzado
con cumplimiento simultáneo de los cuatro umbrales cuantitativos declarados}
(F1~$\geq$~0,80; MAE~$\leq$~0,15; exactitud semántica~$\geq$~0,70;
corrección y fidelidad~$\geq$~0,70).

\subsection*{Verificación de Hipótesis}

A partir de los resultados operacionales medidos sobre el corpus electoral
peruano de marzo de 2026 y de las métricas de calidad obtenidas sobre los
dos \textit{gold datasets} construidos para esta tesis —100 documentos
anotados por humano para el módulo de etiquetado y 30 preguntas anotadas
con \textit{ground truth} para la capa GraphRAG—, se contrasta a
continuación la hipótesis general y cada una de las cinco hipótesis
específicas con la evidencia obtenida:

\textbf{Hipótesis General (HG)} —\textit{El uso integrado de modelos de
lenguaje de gran escala (LLMs) y grafos de conocimiento permite construir un
sistema de inteligencia electoral capaz de realizar análisis contextual,
estructural y relacional del discurso mediático digital con calidad
verificable contra ground truth humano y evaluaciones LLM-as-judge sobre
conjuntos de referencia anotados}— \textbf{se valida}. El sistema integrado
Electoral-Peru-2026 articula los tres componentes (LLMs para ABSA y detección
de temas, grafo Neo4j con 723 fragmentos y 27 temas, e integración de cuatro
fuentes digitales) y produce indicadores de percepción mediática consultables
en lenguaje natural con evidencia trazable al corpus original. La validación
contra \textit{ground truth} humano sobre el dataset de etiquetado (n=100)
respalda la calidad del módulo de análisis con F1 macro de 0.856 en
categoría temática y MAE de 0.107 en sentimiento continuo, mientras que la
evaluación LLM-as-judge sobre el \textit{golden set} de consulta (n=30)
reporta una corrección de respuesta de 0.792 para la capa GraphRAG.

\textbf{HE1} —\textit{Una arquitectura de ingesta y procesamiento basada en
Medallion Architecture permite integrar datos electorales multi-fuente
preservando la consistencia semántica, la trazabilidad y la calidad del
corpus, reduciendo además el costo computacional del análisis LLM mediante
una cascada de filtros pre-modelo}— \textbf{se valida}. El pipeline
implementado articula tres capas de procesamiento (Bronze, Silver, Gold)
sobre infraestructura dual PostgreSQL con \texttt{pgvector} y Neo4j 5.20
Community, integrando las cuatro fuentes digitales (noticias, X/Twitter,
YouTube y TikTok) bajo un esquema canónico común. La cascada de filtros
pre-LLM parametrizada desde catálogos YAML versionados
(\texttt{filter\_by\_account}, \texttt{filter\_by\_hashtag} y
\texttt{filter\_by\_title\_keywords}) redujo de 818 a 655 documentos
($\sim$20\%) el volumen pasado al LLM en la corrida del 22 de marzo de 2026,
con un ahorro computacional asociado de aproximadamente 40\% en costo de
\textit{tokens}. La idempotencia del pipeline se verificó mediante tres
corridas consecutivas con conteos constantes en Postgres
(\texttt{documentos}, \texttt{segmentos}, \texttt{analisis\_llm}) y en Neo4j
(\texttt{Documento}, \texttt{Fragmento}, \texttt{Tema}).

\textbf{HE2} —\textit{Los modelos basados en LLMs en modalidad few-shot
ajustados al contexto político peruano alcanzan métricas de calidad
satisfactorias (F1 macro $\geq 0.80$ en clasificación, MAE $\leq 0.15$ en
sentimiento continuo) tanto en el análisis de sentimiento dirigido como en
la detección contextual de tópicos emergentes del discurso mediático
electoral}— \textbf{se valida}. DeepSeek-Chat en modalidad \textit{few-shot}
con 4 ejemplos calibrados, evaluado contra el \textit{gold dataset} de 100
documentos anotados por humano, obtuvo \textbf{F1 macro 0.856} en categoría
temática (exactitud 0.830 sobre 17 clases), \textbf{F1 0.974} en relevancia
geográfica al Perú, \textbf{F1 macro 0.916} en identificación de candidato
protagonista y \textbf{MAE 0.107} con correlación Pearson de 0.892 en el
sentimiento continuo, superando ambos umbrales declarados. En la tarea de
detección dinámica de temas, el enfoque híbrido LLM + \textit{retrieval}
multinivel alcanzó una exactitud semántica de 0.700 en la asignación
documento$\rightarrow$tema sobre el mismo dataset humano. El único
componente con desempeño limitado fue la extracción multi-etiqueta de
actores principales (F1 micro 0.295).

\textbf{HE3} —\textit{El grafo de conocimiento electoral construido mediante
extracción semántica basada en LLMs permite expresar y consultar relaciones
contextuales entre candidatos, temas, medios, fragmentos y períodos del
ecosistema digital electoral en una única query estructural}— \textbf{se
valida}. El grafo Neo4j Community con sus diez tipos de nodos y siete tipos
de relaciones (\texttt{MENCIONA}, \texttt{SOBRE\_TEMA},
\texttt{PUBLICADO\_POR}, entre otras) permitió ejecutar consultas Cypher
que articulan candidatos, temas, medios, fragmentos y períodos en una sola
expresión. Los conteos del subgrafo del día (723 documentos, 723 fragmentos,
27 temas, 13 categorías) confirman la correcta construcción del grafo y la
correspondencia con los datos relacionales en Postgres. Las catorce
herramientas del servidor MCP, todas implementadas sobre \texttt{MATCH},
\texttt{MERGE} y proyecciones Cypher contra este grafo, demuestran
operativamente que las relaciones modeladas son consultables en producción.

\textbf{HE4} —\textit{La integración híbrida GraphRAG basada en herramientas
tipadas sobre el protocolo MCP mejora la relevancia y la corrección de
respuestas respecto a un sistema RAG por similitud vectorial sobre el mismo
corpus}— \textbf{se valida cuantitativamente}. La métrica directa de
cumplimiento se midió sobre un \textit{golden set} de 30 preguntas anotadas
con \textit{ground truth} de herramientas esperadas y respuesta de
referencia: el LLM cliente (Claude Sonnet) seleccionó el \textbf{conjunto
completo de herramientas esperadas en el 83.3\%} de las preguntas y al
menos una herramienta esperada en el \textbf{100\%}. La evaluación de la
respuesta final mediante un juez independiente (DeepSeek-Chat, $T=0.2$)
reportó una \textbf{corrección de respuesta de 0.792 para el MCP frente a
0.196 para el RAG \textit{baseline}} por similitud vectorial
($\Delta=+0.596$), una \textbf{relevancia de 0.842 vs 0.388}
($\Delta=+0.454$) y una \textbf{similitud de \textit{embeddings} de 0.721
vs 0.579} ($\Delta=+0.142$) sobre el mismo corpus. La ventaja del MCP es
más pronunciada en preguntas \texttt{factual} ($\Delta=+0.902$ en
corrección) y \texttt{causal} ($\Delta=+0.480$), confirmando la
superioridad estructural del paradigma de acceso a datos tabulares por
herramientas tipadas sobre la recuperación textual por \textit{embeddings}.

\textbf{HE5} —\textit{El sistema integrado basado en LLMs, grafos de
conocimiento y GraphRAG alcanza niveles de calidad analítica y coherencia
temática verificables mediante métricas contra ground truth humano y
evaluaciones LLM-as-judge sobre conjuntos anotados}— \textbf{se valida}.
El sistema integrado, operado en producción durante la prueba piloto del
12 y 13 de marzo de 2026, fue evaluado en sus dos planos de salida: (i)
el módulo de etiquetado mostró \textbf{F1 macro 0.856} en categoría temática
y \textbf{MAE 0.107} en sentimiento continuo contra etiquetas humanas
(n=100); (ii) la capa GraphRAG mostró \textbf{tool selection accuracy de
83.3\%}, \textbf{corrección de respuesta de 0.792}, \textbf{relevancia de
0.842} y \textbf{fidelidad de 0.854} sobre un \textit{golden set} de 30
preguntas evaluado por DeepSeek-Chat como juez independiente. La
articulación de ambos planos —etiquetado y consulta— sobre el grafo Neo4j
permite que las consultas en lenguaje natural devuelvan respuestas
sustentadas por evidencia textual del corpus original, cerrando el ciclo
completo del sistema de inteligencia electoral.

\subsection*{Consideraciones y limitaciones}
Sin perjuicio de los resultados obtenidos, el sistema Electoral-Peru-2026
presenta limitaciones metodológicas que deben ser reconocidas. En primer
lugar, la representatividad del universo digital analizado no es equivalente
al electorado real: la población activa en redes sociales presenta un sesgo
hacia grupos etarios más jóvenes y urbanos. En segundo lugar, los modelos
LLM utilizados fueron pre-entrenados principalmente sobre corpus en inglés,
lo que puede introducir sesgos sistemáticos en la clasificación del
sentimiento hacia determinados actores políticos en español. En tercer
lugar, la cobertura del
sistema se limitó a las cuatro fuentes digitales mencionadas, sin incluir
medios audiovisuales tradicionales (radio y televisión) ni plataformas de
mensajería privada como WhatsApp o Telegram, donde circula información que
no es accesible mediante \textit{scraping} público.

En cuarto lugar, los \textit{gold datasets} construidos para la evaluación
contra \textit{ground truth} tienen restricciones de tamaño y diseño que
deben considerarse al interpretar los resultados, y que constituyen
\textbf{limitaciones de validez interna} de la presente tesis: (i) el
dataset de etiquetado (n=100 documentos) se construyó mediante un protocolo
de \textbf{validación asistida por dos evaluadores humanos} que revisaron
de forma conjunta cada documento y resolvieron por consenso las
discrepancias detectadas; dado que el protocolo se basó en consenso
colaborativo y no en anotación independiente paralela, no se reporta el
coeficiente $\kappa$ de Cohen inter-anotador calculado sobre etiquetas
independientes, lo que limita el grado de cuantificación de la fiabilidad
del \textit{ground truth} obtenido; (ii) el
dataset de deduplicación de temas (n=15 pares) tiene una muestra
insuficiente para extraer conclusiones robustas sobre precision/recall de
la fase de \texttt{llm\_dedup}, donde cada caso pesa cerca de 50 puntos
porcentuales; (iii) el \textit{golden set} de consulta (n=30 preguntas)
cubre un período mediático acotado al 10--12 de marzo de 2026, lo que
coloca al RAG \textit{baseline} en desventaja estructural para preguntas
sobre encuestas posteriores a esa fecha.
\section{Recomendaciones}
A partir de los resultados obtenidos y las limitaciones identificadas, se
formulan las siguientes recomendaciones para trabajos de investigación
futuros:

\begin{enumerate}[leftmargin=*,label=\textbf{\arabic*.}]
    \item \textbf{Incorporar medios audiovisuales tradicionales como quinta
    y sexta fuente de datos.} Si bien el sistema Electoral-Peru-2026 ya
    integra las cuatro fuentes digitales más representativas del ecosistema
    electoral peruano (noticias web, X/Twitter, YouTube y TikTok), un
    porcentaje no despreciable del electorado —en particular en zonas
    rurales y entre adultos mayores— consume información política
    principalmente por radio y televisión. Se recomienda incorporar
    transcripciones automáticas de programas de radio (mediante
    \textit{streams} públicos y \texttt{whisper} u otros modelos ASR) y de
    noticieros televisivos (YouTube oficial de los canales) como fuentes
    adicionales del pipeline, lo que permitiría mejorar la representatividad
    del corpus y capturar dinámicas de opinión que hoy no se reflejan en el
    ecosistema digital monitorizado.

    \item \textbf{Realizar ajuste fino (fine-tuning) con datos anotados del
    español político peruano.} Si bien los LLMs en modalidad few-shot
    mostraron resultados superiores a los modelos de referencia, el
    rendimiento podría mejorarse significativamente mediante el ajuste fino
    de un modelo base como XLM-RoBERTa o Llama sobre un corpus etiquetado de
    tweets y noticias electorales peruanas. Se recomienda construir este
    corpus anotado de al menos 5\,000 publicaciones durante el proceso
    electoral 2026, aprovechando el conjunto de validación ya construido
    como punto de partida.

    \item \textbf{Extender el sistema a elecciones subnacionales.} La
    arquitectura de Electoral-Peru-2026 es agnóstica respecto al nivel de la
    contienda electoral. Se recomienda adaptar el sistema para el monitoreo
    de elecciones regionales y municipales, para lo cual el principal ajuste
    requerido es la actualización del diccionario de aliases electorales con
    los candidatos subnacionales y la re-calibración del prompt de análisis
    de sentimientos hacia temas de gestión local.

    \item \textbf{Integrar datos oficiales del JNE y la ONPE.} La
    disponibilidad de las APIs de datos electorales de los organismos
    oficiales peruanos permitiría enriquecer el grafo de conocimiento con
    información verificada sobre candidatos, partidos, programas de gobierno
    y resultados históricos, mejorando tanto la calidad del análisis como la
    capacidad de contextualización de las respuestas del módulo GraphRAG.

    \item \textbf{Definir un protocolo de umbral adaptativo para el ISN.} El
    umbral fijo de $\pm 0.20$ utilizado para la clasificación de sentimientos
    positivo/negativo/neutro fue definido de manera experimental. Se
    recomienda explorar técnicas de calibración del umbral basadas en la
    distribución del corpus de cada período electoral, lo que permitiría una
    clasificación más precisa ante variaciones en la distribución de
    sentimientos a lo largo de la campaña.
\end{enumerate}